\chapter*{Úvod} % chapter* je necislovana kapitola
\addcontentsline{toc}{chapter}{Úvod} % rucne pridanie do obsahu
\markboth{Úvod}{Úvod} % vyriesenie hlaviciek

V posledných desaťročiach zaznamenal svet technológií a informatiky mimoriadny pokrok, pričom jednou z najdôležitejších oblastí tohto vývoja je umelá inteligencia. Umelá inteligencia, ktorá zahŕňa strojové učenie, neurónové siete a autonómne riadenie, má obrovský potenciál zmeniť spôsob, akým žijeme a pracujeme. Od automatizácie priemyselných procesov až po vývoj autonómnych vozidiel sa umelá inteligencia stále viac integruje do našich každodenných životov.

Táto práca sa zaoberá implementáciou autonómneho agenta pre hru Trackmania, simulátor automobilových pretekov, ktorý poskytuje ideálne prostredie na trénovanie a testovanie pokročilých algoritmov umelej inteligencie. Cieľom práce je vytvoriť agenta, ktorý dokáže dokončiť pretekársku trať bez ľudského zásahu, pričom využíva metódy učenia posilňovaním a raycastingu na simuláciu senzorických vnemov.

V teoretickej časti práce sa zaoberáme základmi umelej inteligencie, strojového učenia, neurónových sietí a autonómneho riadenia. Tieto kapitoly poskytujú potrebný kontext a základ pre pochopenie praktických aspektov implementácie agenta. Podrobne vysvetľujeme koncepty ako racionálne správanie, odmeňovacia funkcia, rôzne typy strojového učenia a techniky používané v počítačovej grafike a virtuálnom prostredí.

Praktická časť práce sa zameriava na konkrétnu implementáciu autonómneho agenta v prostredí hry Trackmania. Na rozdiel od existujúcich riešení, ktoré často vyžadujú ľudskú interakciu, spracovanie obrazu alebo 2D aproximácie trate, náš prístup využíva 3D modelovanie a raycasting na získanie presných a realistických senzorických dát. Využívame Python a open-source platformu OpenPlanet na extrahovanie dát z hry a vytváranie virtuálneho prostredia pre trénovanie agenta pomocou učenia posilňovaním.

Výsledky práce ukazujú, že náš agent je schopný úspešne dokončiť pretekárske trate, hoci nedosahuje úroveň skúsených hráčov. V závere práce sumarizujeme dosiahnuté výsledky a navrhujeme možné smery pre zlepšenie nášho autonómneho agenta.

